% 第1章 绪论
\section{绪论}

\subsection{研究背景}

\subsubsection{风电功率预测的现实意义}
随着可再生能源装机规模持续增长，风力发电已成为现代电力系统中的重要组成部分，风电预测也逐渐成为新能源并网运行中的关键研究问题\cite{liu_comprehensive_2025,__2019}。近年来，我国风电装机容量持续快速增长，风电在电网中的渗透率不断提高。然而，风电输出同时受到风速、风向、温度等多种气象因素影响，具有明显的随机性、波动性和间歇性，使其难以像传统火电机组一样实现精确调节，并对电网调峰、调频及经济调度提出更高要求\cite{yang_survey_2024,__2017-1}。

准确的风电功率预测能够为电网调度、备用容量配置及运行风险控制提供可靠依据，合理的误差评价体系也有助于降低系统运行成本并提高电网对风电的消纳能力\cite{__2021,__2011}。在风电功率预测时间尺度上，通常分为超短期预测、短期预测和中长期预测，各尺度的定义与应用场景如表~\ref{tab:wind_power_prediction_time_scale} 所示。

\begin{table}[h]
\centering
\caption{风电功率预测时间尺度与应用场景}
\label{tab:wind_power_prediction_time_scale}
\begin{tabular}{ccc}
\toprule
预测时间尺度 & 预测时间范围 & 主要应用场景 \\
\midrule
超短期预测 & 15分钟$\sim$4小时 & 实时调度、负荷跟踪、储能控制 \\
短期预测 & 4小时$\sim$72小时 & 日前调度计划、电力市场交易 \\
中长期预测 & 3天$\sim$7天及以上 & 电网规划、检修安排 \\
\bottomrule
\end{tabular}
\end{table}

本文重点关注超短期风电功率预测，预测时间范围为未来15分钟至1.5小时（预测步长LF=1至LF=6，每步15分钟）。



\subsection{国内外研究现状}
\subsubsection{国内外研究脉络梳理}
风力发电功率预测研究经历了三个演进阶段。
早期以基于数值天气预报（NWP）的物理模型方法和以 AR/ARIMA 为代表的统计方法为主。
前者重点在于构建风力发电的准确物理模型，后者靠历史功率数和一些统计分析技术进行预测。\cite{liu_deterministic_2019}。

此后，随着神经网络的兴起，ANN、SVM、ELM 等浅层神经网络相继引入，改善了对非线性映射关系的处理。但面对风电功率序列的长期依赖和剧烈波动，浅层神经网络预测能力有限。

2010年中后期，深度学习兴起，LSTM、GRU、CNN等深度学习模型逐渐被引入风电功率预测领域，并进一步向 Attention 机制、Transformer 及图神经网络延伸
\cite{wu_comprehensive_2022, liu_comprehensive_2025, yang_survey_2024}。


\subsubsection{国内研究现状}

国内风电功率预测研究起步于 21 世纪初，早期主要探索"风速—功率"的间接预测路径。
杨秀媛等学者\cite{__2005}提出时序神经网络法，以时间序列模型的基本参数辅助神经网络输入筛选，并指出功率预测难度高于风速预测；
王丽婕等学者\cite{_-_2007}从非线性动力学角度证明风电功率时间序列具有混沌属性，
随后有学者将小波分析与神经网络结合\cite{__2009}，
使预测绝对平均误差由 6.99\% 降至 6.01\%。

2010 年前后，研究重点转向样本构造与泛化能力提升。
孟洋洋等学者\cite{__2010}提出相似日结合 Chebyshev 神经网络的预测模型，
改善了传统 BP 网络收敛慢的问题。
刘克文等学者\cite{__2013}引入集成学习框架，
通过动态调整样本权重提高了模型在突变点的泛化能力。
徐曼等学者\cite{__2011}提出涵盖纵向、横向与极端误差的综合评价体系，
推动国内研究从单一均值指标转向多维评价。

2017 年以来，深度学习成为主流方向。
黎静华等学者\cite{__2017-1}系统梳理了点预测、区间预测、概率预测与场景预测四条主线；
朱乔木等学者\cite{__2017}率先将多变量 LSTM 引入超短期预测，
验证了其相对于 ANN 和 SVM 的时序建模优势；
韩自奋等学者\cite{__2019}进一步指出空间相关性预测与爬坡预测是重要新兴方向；
孙荣富等学者\cite{__2021}从 NWP 到调度应用的全流程视角归纳了工程落地路径。
谢小瑜等学者\cite{_w-bilstm_2021}构建小波分解结合双向 LSTM 与注意力机制的复合框架；
李大中等学者\cite{__2021-1}提出 BiGRU 叠加随机森林误差修正与核密度估计的一体化框架，
体现出"深度网络 + 误差修正 + 区间预测"的研究范式逐步成熟。

近两三年，国内研究呈现出混合架构持续深化与研究目标多元化并进的趋势。
陈大为等学者\cite{_lstmxgboost_2025}提出 LSTM 与 XGBoost 的混合模型，
由 LSTM 提取时序依赖特征后再由 XGBoost 进行非线性拟合修正，
在预测精度和拟合能力上均优于单一 LSTM；
王丽杰等学者\cite{__2022-1}耦合卷积层、GRU 与全连接网络，
基于 SCADA 数据开展风电机组功率预测，并将该框架延伸至性能劣化检测；
韩顺天学者\cite{__2025-2}将 VMD 分解与 Informer 编码结合，
通过频域增强通道注意力机制实现多频特征融合；
张东学者\cite{__2025-1}则结合模糊信息粒化与时间卷积网络提取局部特征，
再输入改进 Transformer 进行超短期预测，区间覆盖率达到 0.944。
马昕远和伊海港学者\cite{_cnn-lstm-attention_2025}构建 CNN-LSTM-Attention 模型，
并引入 SHAP 方法对湿度、风速、气压等特征贡献进行量化解释。


\subsection{本文主要研究内容}

针对风电功率序列随机性强、多步预测误差累积等问题。本文以六个实际风电场的历史运行数据和相关气象数据为基础，使用CNN、LSTM和CNN-LSTM三种深度学习模型对风电功率进行超短期预测，主要研究内容如下：

（1）风电功率预测数据集的数据分析和预处理。针对原始风电数据中存在的缺失值、异常值、量纲不一致等问题，采用插值、IQR异常值处理、标准化等方法对数据进行清洗。并将时间这一列数据进行正余弦编码变成新特征输入。最后采用滑动时间窗口的方法把时间序列构造为深度学习所需的输入输出张量。最后进行数据集划分。

（2）典型深度学习预测模型的原理分析与建模。介绍CNN、LSTM及CNN-LSTM混合模型的基本原理与网络结构，分析各模型在局部特征提取、时序依赖建模等方面的特点与适用性。最后设计各个模型的超参数和结构。

（3）多模型在多风场、多步长下的对比实验。在六个不同装机容量和地理条件的风电场上，对三种模型进行不同步长的实验。最后通过计算测试集预测的MSE、RMSE、MAE和$R^2$等评价指标，从预测精度、预测步长对不同模型精度的影响，不同风电场的精度等多个维度进行分析。

\subsection{本文结构安排}

本文共分为五章，各章内容安排如下：

第1章为绪论。介绍风电功率预测的研究背景与现实意义，讨论国内外相关研究现状以及说明本文的主要研究内容与各章节安排。

第2章为数据集介绍与预处理。介绍实验数据集的来源与基本特征，详细描述数据清洗、特征工程及滑动窗口样本构造方法，为后续模型训练提供数据基础。

第3章为深度学习预测模型。介绍CNN、LSTM和CNN-LSTM混合模型的基本原理与网络结构，并说明本文采用的模型评价指标和实验环境设置。

第4章为算例分析。在六个风电场上对三种模型进行超短期预测实验，从预测精度、预测步长对不同模型精度的影响，不同风电场的精度等多个维度进行分析。s

第5章为结论与展望。总结本文的主要研究结论，并对未来可能的研究方向进行展望。
